% Created 2026-05-07 qui 08:30
% Intended LaTeX compiler: xelatex
\documentclass[12pt,a4paper]{article}
\usepackage{graphicx}
\usepackage{longtable}
\usepackage{wrapfig}
\usepackage{rotating}
\usepackage[normalem]{ulem}
\usepackage{capt-of}
\usepackage{hyperref}
\usepackage[T1]{fontenc}
\usepackage[utf8]{inputenc}
\usepackage[brazil, brazilian]{babel}
\usepackage[top=2.5cm,bottom=2.5cm,left=3cm,right=2.5cm]{geometry}
\usepackage{lmodern}
\usepackage{microtype}
\usepackage{minted}
\setminted{fontsize=\small, linenos=false, breaklines=true, bgcolor=gray!6}
\usepackage[colorlinks=true, linkcolor=black, urlcolor=blue, citecolor=black]{hyperref}
\usepackage{fancyhdr}
\pagestyle{fancy}
\fancyhf{}
\fancyhead[L]{\small INF01090 — Ciência de Dados}
\fancyhead[R]{\small Laboratório 5}
\fancyfoot[C]{\thepage}
\renewcommand{\headrulewidth}{0.4pt}
\usepackage{setspace}
\setstretch{1.15}
\usepackage{fancyvrb}
\usepackage{etoolbox}
\AtBeginDocument{\setlength{\parskip}{4pt}}
\author{Nícolas Rosenthal Dal Corso}
\date{Ciência de Dados (INF 01090/2026) \\\\ May, 2026}
\title{Laboratório 5 \\ Relatório}
\hypersetup{
 pdfauthor={Nícolas Rosenthal Dal Corso},
 pdftitle={Laboratório 5 \\ Relatório},
 pdfkeywords={},
 pdfsubject={},
 pdfcreator={Emacs 30.2 (Org mode 9.7.11)}, 
 pdflang={Portuges}}
\begin{document}

\maketitle
\setcounter{tocdepth}{3}
\tableofcontents

\section{Visão Geral}
\label{sec:org740da2d}

Este relatório descreve a implementação do \emph{pipeline} de predição de preços de
imóveis desenvolvido para a competição \emph{House Prices - Advanced Regression
Techniques} (Kaggle), adaptada para o conjunto de dados de avaliação da
disciplina \texttt{INF01090/2026}.

O modelo final é um \textbf{stack ensemble} de dois níveis com meta-blending via
Optuna, cujo OOF RMSE (em escala log) obtido foi de \textbf{0.112577}.

Todas as submissões ocorreram sob o nome \texttt{n-rosenthal}.
\section{Dados e Pré-processamento}
\label{sec:orgba46b6b}

\subsection{Conjunto de Dados}
\label{sec:org8b6382b}

\begin{center}
\begin{tabular}{ll}
Atributo & Valor\\
\hline
Linhas (treino) & 1 016 (após remoção de outliers)\\
Linhas (teste) & 438\\
Features brutas & 80 variáveis\\
Features finais & 125 (após engenharia)\\
Target & \texttt{SalePrice} (numérica, contínua)\\
\end{tabular}
\end{center}
\subsection{Remoção de Outliers}
\label{sec:orgcb8dd11}

Cinco regras de filtragem foram aplicadas sobre o conjunto de treino:

\begin{verbatim}
train_df = train_df[train_df["GrLivArea"] < 4500]
train_df = train_df[~((train_df["GrLivArea"] > 4000) & (train_df["SalePrice"] < 200_000))]
train_df = train_df[train_df["LotArea"] < 100_000]
train_df = train_df[train_df["TotalBsmtSF"] < 6_000]
train_df = train_df[~((train_df["OverallQual"] <= 4) & (train_df["SalePrice"] > 300_000))]
\end{verbatim}

Essas regras removem imóveis com \emph{living area} muito grande, lotes
exageradamente grandes e combinações de baixa qualidade com preço alto, que
distorceriam o aprendizado.
\subsection{Transformação da Target}
\label{sec:orgdc91d38}

O preço de venda (\texttt{SalePrice}) foi transformado via \texttt{log1p} antes do
treinamento e revertido via \texttt{expm1} nas predições finais. Isso reduz a
assimetria da distribuição e estabiliza o treinamento dos modelos de boosting.
\subsection{Correção de Tipos}
\label{sec:org1667ff4}

\texttt{MSSubClass}, \texttt{MoSold} e \texttt{YrSold} são variáveis \textbf{nominais}, não números
contínuos. A conversão para \texttt{string} é feita \textbf{antes} da separação entre
\texttt{cat\_cols} e \texttt{num\_cols}, garantindo que sejam tratadas como categorias pelo
\texttt{TargetEncoder} e não como grandezas ordenadas.

\begin{verbatim}
def fix_dtypes(df):
    for col in ["MSSubClass", "MoSold", "YrSold"]:
        if col in df.columns:
            df[col] = df[col].astype(str)
    return df
\end{verbatim}
\subsection{Imputação de Valores Ausentes}
\label{sec:org40d6b3c}

Duas estratégias foram adotadas:

\begin{itemize}
\item \textbf{Colunas categóricas sem valor} (\texttt{Alley}, \texttt{MasVnrType}, \texttt{MiscFeature}):
preenchidas com a string \texttt{"None"}, que é mapeada para 0 nos encodings
ordinais.

\item \textbf{Colunas numéricas} (\texttt{MasVnrArea}, \texttt{BsmtFinSF1/2}, \texttt{GarageArea}, etc.):
preenchidas com \texttt{0}, representando ausência do atributo.

\item \textbf{\texttt{GarageYrBlt}}: imputado com \texttt{YearBuilt} do próprio imóvel (não com \texttt{0},
que geraria um ano inválido e introduziria viés semântico).
\end{itemize}

\begin{verbatim}
mask = df["GarageYrBlt"].isna()
df.loc[mask, "GarageYrBlt"] = df.loc[mask, "YearBuilt"]
\end{verbatim}
\subsection{Encoding Ordinal}
\label{sec:org6cd0c2c}

Variáveis de qualidade com escala natural (\texttt{Ex > Gd > TA > Fa > Po}) foram
codificadas numericamente por mapeamento direto, evitando que o modelo trate
categorias ordenadas como nominais:

\begin{center}
\begin{tabular}{ll}
Coluna & Escala\\
\hline
\texttt{ExterQual} & Ex=5, Gd=4, TA=3, Fa=2, Po=1, NA=0\\
\texttt{BsmtQual} & (idem acima)\\
\texttt{KitchenQual} & (idem acima)\\
\texttt{BsmtFinType1} & GLQ=6, ALQ=5, … Unf=1, NA=0\\
\texttt{BsmtExposure} & Gd=4, Av=3, Mn=2, No=1, NA=0\\
\texttt{Functional} & Typ=8, Min1=7, … Sal=1\\
\texttt{GarageFinish} & Fin=3, RFn=2, Unf=1, NA=0\\
\end{tabular}
\end{center}
\section{Engenharia de Features}
\label{sec:org4f9480b}

Foram criadas 45 features adicionais, organizadas em sete grupos:
\subsection{Área Total e Contagens}
\label{sec:org101d0ee}

\begin{verbatim}
df["TotalSF"]        = df["TotalBsmtSF"] + df["1stFlrSF"] + df["2ndFlrSF"]
df["TotalBathrooms"] = df["FullBath"] + 0.5*df["HalfBath"] + \
                       df["BsmtFullBath"] + 0.5*df["BsmtHalfBath"]
df["TotalPorchSF"]   = df["OpenPorchSF"] + df["EnclosedPorch"] + \
                       df["3SsnPorch"]   + df["ScreenPorch"]
df["GarageScore"]    = df["GarageCars"] * df["GarageArea"]
\end{verbatim}
\subsection{Features Temporais}
\label{sec:org577b28c}

\begin{verbatim}
df["HouseAge"]       = df["YrSold"].astype(int) - df["YearBuilt"]
df["RemodAge"]       = df["YrSold"].astype(int) - df["YearRemodAdd"]
df["TimeSinceRemod"] = df["YrSold"].astype(int) - df["YearRemodAdd"]
df["IsNew"]          = (df["YrSold"].astype(int) == df["YearBuilt"]).astype(int)
df["IsRemodeled"]    = (df["YearRemodAdd"] != df["YearBuilt"]).astype(int)
\end{verbatim}
\subsection{Flags Binárias}
\label{sec:orgf539038}

Nove flags indicam presença/ausência de atributos importantes:
\texttt{HasPool}, \texttt{Has2ndFloor}, \texttt{HasGarage}, \texttt{HasBsmt}, \texttt{HasFireplace},
\texttt{IsRemodeled}, \texttt{IsNew}, \texttt{HasMasVnr}, \texttt{HasPorch}.
\subsection{Interações Multiplicativas}
\label{sec:orgad435f1}

\begin{verbatim}
df["OverallQual_TotalSF"]   = df["OverallQual"] * df["TotalSF"]
df["OverallQual_GrLivArea"] = df["OverallQual"] * df["GrLivArea"]
df["OverallQual_sq"]        = df["OverallQual"] ** 2
df["QualCondScore"]         = df["OverallQual"] * df["OverallCond"]
df["QualAge"]               = df["OverallQual"] / df["HouseAge"].clip(1)
df["NewAndQual"]            = df["IsNew"] * df["OverallQual"]
\end{verbatim}
\subsection{Log-Transforms de Features Assimétricas}
\label{sec:org16fb799}

\begin{verbatim}
df["LotAreaLog"]    = np.log1p(df["LotArea"])
df["GrLivAreaLog"]  = np.log1p(df["GrLivArea"])
df["TotalSFLog"]    = np.log1p(df["TotalSF"])
df["GarageAreaLog"] = np.log1p(df["GarageArea"])
\end{verbatim}
\subsection{Razões (\emph{Ratios})}
\label{sec:org69fdc3c}

\begin{verbatim}
df["LivAreaRatio"]   = df["GrLivArea"]   / df["TotalSF"].clip(1)
df["BsmtRatio"]      = df["TotalBsmtSF"] / df["TotalSF"].clip(1)
df["GarageRatio"]    = df["GarageArea"]  / df["GrLivArea"].clip(1)
df["LotFrontRatio"]  = df["LotFrontage"].fillna(0) / df["LotArea"].clip(1)
df["BsmtFinRatio"]   = df["BsmtFinSF1"]  / df["TotalBsmtSF"].clip(1)
df["LotAreaPerRoom"] = df["LotArea"]     / df["TotRmsAbvGrd"].clip(1)
df["SFPerRoom"]      = df["GrLivArea"]   / df["TotRmsAbvGrd"].clip(1)
\end{verbatim}
\subsection{Scores Compostos}
\label{sec:orgafdcb3b}

\begin{verbatim}
df["ExterScore"]      = df["ExterQual"]    * df["ExterCond"]
df["KitchenScore"]    = df["KitchenQual"]  * df["GrLivArea"]
df["FireScore"]       = df["FireplaceQu"]  * df["Fireplaces"]
df["BsmtScore"]       = df["BsmtExposure"] * df["BsmtQual"]
df["GarageFullScore"] = df["GarageQual"]   * df["GarageCars"]
\end{verbatim}
\subsection{Sazonalidade de Venda}
\label{sec:orgca01b0e}

\begin{verbatim}
mo = df["MoSold"].astype(int)
df["SaleSeasonQ1"] = mo.isin([1, 2, 3]).astype(int)
# Q2, Q3, Q4 — idem
\end{verbatim}
\section{Arquitetura do Modelo}
\label{sec:orgb6dca74}

O \emph{pipeline} adota uma arquitetura de \textbf{stack ensemble} em dois níveis.
\subsection{Nível 1 — Base Learners}
\label{sec:org58ef214}

Oito modelos de gradient boosting formam o primeiro nível:

\begin{center}
\begin{tabular}{llrl}
ID & Algoritmo & Profundidade & Observações\\
\hline
\texttt{cat\_a} & CatBoost & 5 & Config principal; usa Optuna\\
\texttt{cat\_b} & CatBoost & 6 & Mais árvores, LR menor\\
\texttt{cat\_c} & CatBoost & 4 & Máxima regularização\\
\texttt{cat\_d} & CatBoost & 7 & Lado mais profundo\\
\texttt{cat\_e} & CatBoost (Lossguide) & 6 & Política de crescimento diferente\\
\texttt{lgb\_best} & LightGBM & 5 & \texttt{num\_leaves=20}; usa Optuna\\
\texttt{xgb\_best} & XGBoost & 4 & \texttt{gamma=0.1}, conservador\\
\texttt{hgb} & HistGradBoost (sklearn) & — & Implementação ortogonal\\
\end{tabular}
\end{center}

Os modelos \texttt{CatBoost} dominaram o ranking OOF (0.114–0.117 RMSE), enquanto
LGB, XGB e HGB ficaram na faixa 0.124–0.126, contribuindo com diversidade ao
ensemble.
\subsection{Busca de Hiperparâmetros com Optuna}
\label{sec:orgd709ed2}

Os hiperparâmetros de \texttt{cat\_a} e \texttt{lgb\_best} são encontrados por busca
bayesiana com \texttt{Optuna} (\texttt{n\_trials=50}), usando OOF RMSE em validação cruzada
de 5 folds como métrica. O espaço de busca cobre \texttt{iterations},
\texttt{learning\_rate}, \texttt{depth}, \texttt{l2\_leaf\_reg}, \texttt{bagging\_temperature},
\texttt{border\_count} e \texttt{random\_strength} (CatBoost), e os parâmetros análogos do
LGB. Os parâmetros dos demais modelos (\texttt{cat\_b} a \texttt{cat\_e}, \texttt{xgb\_best}, \texttt{hgb})
são fixos e cobrem regiões complementares do espaço de hipóteses, garantindo
diversidade.
\subsection{Validação Cruzada (10-fold, 3 repetições)}
\label{sec:orge46a8cb}

O treinamento dos base learners utiliza \texttt{RepeatedKFold} com 10 splits e
\textbf{3 repetições} (\texttt{N\_REPEATS=3}, totalizando 30 folds), o que reduz a
variância da estimativa OOF em \texttt{√3 ≈ 1,73×}.

Para cada fold:

\begin{enumerate}
\item \texttt{SimpleImputer(median)} + \texttt{RobustScaler} aplicados \textbf{apenas nos dados de treino}.
\item \texttt{TargetEncoder(smoothing=10)} ajustado \textbf{apenas nos dados de treino}.
\item \texttt{PCA(n\_components=10)} sobre o bloco categórico \emph{encoded}.
\item Concatenação: \texttt{[numérico\_escalado | cat\_encoded | PCA\_cat]}.
\item Treinamento com early stopping (300 rounds para \texttt{CatBoost}, 200 para \texttt{LGB/XGB}).
\end{enumerate}

Este procedimento garante \textbf{zero data leakage} entre folds.
\subsection{Nível 2 — Meta-Learners}
\label{sec:org30e4ae5}

Quatro estratégias de blending são treinadas sobre as predições OOF,
enriquecidas com \textbf{features de 2ª ordem} (média, desvio-padrão, máximo, mínimo
e amplitude entre os 8 modelos):

\begin{center}
\begin{tabular}{lrl}
Meta-Learner & OOF RMSE & Notas\\
\hline
\texttt{RidgeCV} & 0.112582 & Treinado sobre features de 2ª ordem\\
\texttt{ElasticNetCV} & 0.112599 & \$\(\alpha\)\$=0.00007, \$l\_1\$=0.90\\
\texttt{LassoCV (+)} & 0.113698 & \texttt{positive=True}; pesos: cat\_c=0.35, cat\_e=0.68\\
Optuna (direto) & 0.114212 & Busca \emph{simplex} sobre OOF base\\
\end{tabular}
\end{center}
\subsection{Meta-Blending Final (Optuna)}
\label{sec:org42197cc}

Um segundo nível de Optuna combina os quatro meta-learners:

\begin{verbatim}
★ Pesos finais:
  ridge        : 0.6167
  elasticnet   : 0.3831
  lasso_pos    : 0.0001
  direct_opt   : 0.0001

★ OOF RMSE FINAL: 0.112577
\end{verbatim}
\subsection{Pseudo-Labelling (2ª Passagem)}
\label{sec:org9f3736d}

Após o blending final, amostras de teste com baixa variância entre as
predições dos 8 modelos (percentil 30 do desvio-padrão) são adicionadas ao
conjunto de treino com suas labels previstas. O \emph{pipeline} completo é
re-treinado sobre esse conjunto aumentado, e as predições finais são extraídas
da 2ª passagem.

\begin{verbatim}
std_test  = test_matrix.std(axis=1)
confident = std_test < np.percentile(std_test, 30)
X_aug = pd.concat([X, X_test[confident]], ignore_index=True)
y_aug = pd.concat([y, pd.Series(final_log[confident])], ignore_index=True)
# → re-treino completo do CV loop sobre X_aug, y_aug
\end{verbatim}
\section{Pré-processamento \texttt{Inside-Fold}}
\label{sec:org752206c}

\begin{verbatim}
Para cada fold k:
  ┌─────────────────────────────────────────────────────────────────┐
  │ X_train_k, X_val_k = split(X, fold=k)                           │
  │                                                                 │
  │ imputer.fit(X_train_k[num])   → X_train_num, X_val_num          │
  │ scaler.fit(X_train_k[num])    → (RobustScaler)                  │
  │                                                                 │
  │ TargetEncoder.fit(X_train_k, y_train_k)                         │
  │   → X_train_te, X_val_te, X_test_te                             │
  │                                                                 │
  │ PCA(10).fit(X_train_te)       → X_train_pca, X_val_pca          │
  │                                                                 │
  │ X_full = hstack[num | te | pca]                                 │
  │                                                                 │
  │ Para cada base model m:                                         │
  │   model.fit(X_train_full, y_train, eval=(X_val_full, y_val))    │
  │   oof[val_idx, m] += model.predict(X_val_full)                  │
  │   test_preds[m]   += model.predict(X_test_full) / n_folds       │
  └─────────────────────────────────────────────────────────────────┘
\end{verbatim}
\section{Resultados}
\label{sec:org7ea2461}

\subsection{RMSE por Base Learner (OOF, escala logarítmica)}
\label{sec:orgf807026}

\begin{center}
\begin{tabular}{lr}
Modelo & OOF RMSE\\
\hline
\texttt{cat\_e} & 0.114478\\
\texttt{cat\_c} & 0.115069\\
\texttt{cat\_a} & 0.115290\\
\texttt{cat\_b} & 0.116325\\
\texttt{cat\_d} & 0.117181\\
\texttt{hgb} & 0.124021\\
\texttt{lgb\_best} & 0.124906\\
\texttt{xgb\_best} & 0.125658\\
\end{tabular}
\end{center}
\subsection{Resultado Final}
\label{sec:org4c007f7}

\begin{center}
\begin{tabular}{ll}
Métrica & Valor\\
\hline
OOF RMSE (log) & \textbf{0.112577}\\
Escala da predição & Preço em USD\\
Range de predições & {[}46 419 , 569 716]\\
Arquivo de submissão & \texttt{n-rosenthal-{}-{}submission-6.csv}\\
\end{tabular}
\end{center}
\section{Conclusão}
\label{sec:org3cdf2a8}

O \emph{pipeline} implementado alcançou OOF RMSE de \textbf{0.112577} em escala log,
correspondendo a um erro médio de aproximadamente ±12\% no preço de venda.

Os principais fatores de sucesso foram:

\begin{enumerate}
\item \textbf{Engenharia de features cuidadosa}: 45 features derivadas de conhecimento
do domínio imobiliário, com destaque para as interações \texttt{OverallQual × TotalSF}
e os novos scores compostos (\texttt{ExterScore}, \texttt{FireScore}, \texttt{BsmtScore}).

\item \textbf{Correção de tipos antes da separação de colunas}: converter \texttt{MSSubClass},
\texttt{MoSold} e \texttt{YrSold} para \texttt{string} antes de definir \texttt{cat\_cols\textasciitilde{}/\textasciitilde{}num\_cols}
evita que variáveis nominais sejam tratadas como contínuas.

\item \textbf{Imputação semântica de \texttt{GarageYrBlt}}: usar \texttt{YearBuilt} em vez de \texttt{0}
elimina um viés que confunde os modelos de boosting.

\item \textbf{Target encoding dentro do fold}: eliminou \emph{data leakage} e permitiu
capturar variações de preço por bairro (\texttt{Neighborhood}).

\item \textbf{Diversidade no ensemble}: combinar CatBoost (cinco variantes com
configurações distintas), LGB, XGB e HGB gerou ganho consistente sobre
qualquer modelo individual.

\item \textbf{Meta-blending via Optuna com features de 2ª ordem}: a adição de
estatísticas (média, desvio, range) sobre as predições OOF enriqueceu
o meta-learner; a combinação de Ridge e ElasticNet pesada pelo Optuna
superou qualquer meta-learner isolado.

\item \textbf{Pseudo-labelling com re-treino completo}: a 2ª passagem do CV loop sobre
o conjunto aumentado com amostras de teste de alta confiança contribuiu
para redução adicional do erro.
\end{enumerate}
\end{document}
